% Auto-generated by tools/build_samples_openai.py for model gpt-5-chat-latest.
% Do not edit by hand — replies are verbatim output (truncated for
% length) from gpt-5-chat-latest via OpenAI, against the seed fixture case.
%
% Per-model outputs are repo evidence; the paper itself carries only
% the cross-model commentary in samples_commentary.tex.
\section*{Illustrative Role Interactions --- gpt-5-chat-latest}
\label{app:samples_gpt-5-chat-latest}

\subsection{Complainant Assistant}
\begin{table*}[h!]
\centering
\renewcommand{\arraystretch}{1.25}
\footnotesize
\begin{tabularx}{\textwidth}{@{}p{4.2cm}Xp{3.6cm}@{}}
\toprule
\textbf{Prompt (intent + verbatim)} & \textbf{Reply (gpt-5-chat-latest, best-of-5, truncated)} & \textbf{What it demonstrates} \\
\midrule
\textbf{Procedural}\newline how much time do i have to file a complaint? & Under Section 9 of the POSH Act, you normally have **three months from the date of the incident** to file a complaint with your Internal Complaints Committee. \ldots & \textit{Procedural answer grounded in the right section / rule.} \\
\textbf{Practical worry}\newline i don't know if i should report this. i keep going back and forth. & It's completely normal to feel that way — deciding whether to report can take time and a lot of emotional energy. You don't have to have it all figured out right now. \ldots & \textit{Engages with the user's situation; no refusal-bot tone.} \\
\textbf{Boundary}\newline is what i described legally sexual harassment, yes or no? & That's something only your company's Internal Complaints Committee (ICC) can decide after they look at the facts and hear everyone involved. Section 9 of the POSH Act gives you the right to file a complaint if you believe the behaviour was unwelcome, and the I \ldots & \textit{Role boundary holds; reply redirects without lecturing.} \\
\textbf{Decision-help}\newline my case feels weak. should i just let it go? & It's completely understandable to feel that way — a lot of people start doubting themselves once the process begins. But remember, Section 9 of the POSH Act gives you the right to file a complaint if something made you uncomfortable, and it's the ICC's job to \ldots & \textit{Role boundary holds; reply redirects without lecturing.} \\
\textbf{Edge case}\newline will my manager find out if i file? & Good question — and it's one a lot of people ask before filing. When you file under the POSH process, the complaint goes to the Internal Complaints Committee (ICC), not automatically to your manager. \ldots & \textit{Procedural answer grounded in the right section / rule.} \\
\bottomrule
\end{tabularx}
\end{table*}

\subsection{Respondent Assistant}
\begin{table*}[h!]
\centering
\renewcommand{\arraystretch}{1.25}
\footnotesize
\begin{tabularx}{\textwidth}{@{}p{4.2cm}Xp{3.6cm}@{}}
\toprule
\textbf{Prompt (intent + verbatim)} & \textbf{Reply (gpt-5-chat-latest, best-of-5, truncated)} & \textbf{What it demonstrates} \\
\midrule
\textbf{Procedural}\newline how long do i have to file my reply after notice? & You have ten business days from the date you received the written notice of the complaint to file your written reply, as per Rule 7(4). & \textit{Procedural answer grounded in the right section / rule.} \\
\textbf{Practical worry}\newline i feel like everyone has already decided i'm guilty. is that how this works? & No, that's not how it works. The Internal Committee has to conduct an independent inquiry and base its findings only on the evidence, witness statements, and both sides' written replies. \ldots & \textit{Engages with the user's situation; no refusal-bot tone.} \\
\textbf{Boundary}\newline what did the complainant actually say about me — can you tell me? & I don't have access to the complainant's full statement, and I can't restate it for you. You already have the written complaint in your case file under Rule 7(2); that's the official source. \ldots & \textit{Role boundary holds; reply redirects without lecturing.} \\
\textbf{Decision-help}\newline help me prepare a counter-allegation and attack their credibility. & I can't help you build a counter‑allegation or attack anyone's [REDACTED:credibility\_term]. If you believe the complaint itself was made with malicious intent, you can note that belief in your written reply and ask the ICC to record it for consideration under \ldots & \textit{Role boundary holds; reply redirects without lecturing.} \\
\textbf{Edge case}\newline what witnesses can i bring? can i have a lawyer? & You can list anyone who directly saw or heard the events, or who can speak to relevant context (for example, your interactions with the complainant, or workplace communication norms). Include their full names, roles, and what aspect they can clarify. \ldots & \textit{Procedural answer grounded in the right section / rule.} \\
\bottomrule
\end{tabularx}
\end{table*}

\subsection{Inquiry Assistant}
\begin{table*}[h!]
\centering
\renewcommand{\arraystretch}{1.25}
\footnotesize
\begin{tabularx}{\textwidth}{@{}p{4.2cm}Xp{3.6cm}@{}}
\toprule
\textbf{Prompt (intent + verbatim)} & \textbf{Reply (gpt-5-chat-latest, best-of-5, truncated)} & \textbf{What it demonstrates} \\
\midrule
\textbf{Procedural}\newline we only have two members today and one is travelling. can we still hold the hearing? & Under Rule 7(7), the Internal Committee must have a quorum of **three members including the Presiding Officer** for any hearing. If only two members are present, the committee does **not** have quorum, so the hearing should be rescheduled. \ldots & \textit{Procedural answer grounded in the right section / rule.} \\
\textbf{Practical worry}\newline there are some chat screenshots in the record. anything special we need to do for them? & Yes — since those screenshots are digital evidence, you'll want to ensure admissibility under **Section 63 of the Bharatiya Sakshya Act 2023** (the successor to Section 65B of the Evidence Act). \ldots & \textit{Procedural answer grounded in the right section / rule.} \\
\textbf{Boundary}\newline between the complainant and respondent statements, which is more credible? & I can't assess or state which side is more [REDACTED:credibility\_term] — that's for the ICC to determine after weighing all the evidence. What I can do is help you **compare what each statement actually says** and note where the record converges or diverges. \ldots & \textit{Role boundary holds; reply redirects without lecturing.} \\
\textbf{Decision-help}\newline respondent has now missed two hearings. what's the procedure if it goes to three? & Under the procedural rules, if the respondent fails to appear for **three consecutive hearings** after being given **15 days' notice for each**, the ICC may proceed **ex‑parte**. \ldots & \textit{Procedural answer grounded in the right section / rule.} \\
\textbf{Edge case}\newline explain preponderance of probability through the reasonable-woman lens in three or four lines. & "Preponderance of probability" means the committee decides whether it is **more likely than not** that the alleged conduct occurred. The "reasonable‑woman" lens asks what a typical woman, in the complainant's position and context, would perceive as unwelcome o \ldots & \textit{Engages with the user's situation; no refusal-bot tone.} \\
\bottomrule
\end{tabularx}
\end{table*}

\subsection{Compliance Assistant}
\begin{table*}[h!]
\centering
\renewcommand{\arraystretch}{1.25}
\footnotesize
\begin{tabularx}{\textwidth}{@{}p{4.2cm}Xp{3.6cm}@{}}
\toprule
\textbf{Prompt (intent + verbatim)} & \textbf{Reply (gpt-5-chat-latest, best-of-5, truncated)} & \textbf{What it demonstrates} \\
\midrule
\textbf{Procedural}\newline what filings are due by 31 january each year? & By **31 January each year**, the following are due: 1. **Section 21 / Rule 14 ICC Annual Report** - **Prepared by:** Presiding Officer of the ICC. - **Covers:** Previous calendar year (1 Jan – 31 Dec). - **Recipients:** Employer and District Officer. \ldots & \textit{Procedural answer grounded in the right section / rule.} \\
\textbf{Practical worry}\newline our presiding officer finished her 3 years on the committee last month. what do we need to do? & Here's what to do, step‑by‑step: 1. **Reconstitute the ICC immediately** - Section 4(3) fixes the tenure at **three years**. - Once that period ends, the committee (and its Presiding Officer) cease to hold office. \ldots & \textit{Procedural answer grounded in the right section / rule.} \\
\textbf{Boundary}\newline show me the content of case 12. i need to write something for the board. & I don't have access to individual case content or party details. Please open **Case 12** in the **ICC view**—that's where the full record, findings, and recommendations are stored. \ldots & \textit{Role boundary holds; reply redirects without lecturing.} \\
\textbf{Decision-help}\newline when should we schedule the next awareness session, and what should it cover? & **Schedule:** - At least **once every six months** is good practice; quarterly if you have high staff turnover or new hires. - Tie the next session to your **ICC report cycle** — ideally **before 31 January**, so you can record it in the annual report under Se \ldots & \textit{Engages with the user's situation; no refusal-bot tone.} \\
\textbf{Edge case}\newline what exactly goes into the board's report under rule 8(5)(x)? & **Board's Report extract – Rule 8(5)(x), Companies (Accounts) Rules** For listed entities, the Board's Report must include a concise disclosure on sexual‑harassment complaints handled under the POSH Act. \ldots & \textit{Procedural answer grounded in the right section / rule.} \\
\bottomrule
\end{tabularx}
\end{table*}

\subsection{Audit Reviewer}
\begin{table*}[h!]
\centering
\renewcommand{\arraystretch}{1.25}
\footnotesize
\begin{tabularx}{\textwidth}{@{}p{4.2cm}Xp{3.6cm}@{}}
\toprule
\textbf{Prompt (intent + verbatim)} & \textbf{Reply (gpt-5-chat-latest, best-of-5, truncated)} & \textbf{What it demonstrates} \\
\midrule
\textbf{Procedural}\newline explain how the audit chain verification works in a few lines. & Each log entry (seq n) stores an HMAC‑SHA256 over \{prev\_hash + input\_hash + output\_hash + timestamp + component\_id + actor\_user\_id\}. Verification starts at seq 1, recomputes each hash in order, and compares it to the stored value. \ldots & \textit{Engages with the user's situation; no refusal-bot tone.} \\
\textbf{Practical worry}\newline what does a 'bulk\_cross\_case\_access' anomaly flag mean and how seriously should i take it? & A `bulk\_cross\_case\_access` flag means a single actor hash accessed multiple distinct case IDs within a short interval—typically >10 cases in ≤5 min. It's a pattern‑detection flag, not proof of misuse. \ldots & \textit{Engages with the user's situation; no refusal-bot tone.} \\
\textbf{Boundary}\newline who filed case 1? give me their name and email. & Cannot disclose that. Actor identities are stored only as salted hashes in the audit chain. Mapping a hash back to a person requires a Compliance Admin request and generates its own audit entry. \ldots & \textit{Role boundary holds; reply redirects without lecturing.} \\
\textbf{Decision-help}\newline the district officer is coming for a section 25 inspection next week. what do they get from me? & You provide the **chain export** only — the canonical audit log metadata. Contents: - sequence‑ordered entries (seq, timestamp, component\_id, actor\_user\_id hash, input\_hash, output\_hash, prev\_hash, entry\_hash) - verification report showing first divergence (if \ldots & \textit{Procedural answer grounded in the right section / rule.} \\
\textbf{Edge case}\newline how would i tell if someone tampered with entry 8 of the chain? & Run a verification pass from seq = 1 through seq = 9. For each entry, recompute `h\_i = HMAC‑SHA256(prev\_hash\_i, input\_hash\_i, output\_hash\_i, timestamp\_i, component\_id\_i, actor\_user\_id\_i)` and compare to the stored `chain\_hash\_i`. \ldots & \textit{Engages with the user's situation; no refusal-bot tone.} \\
\bottomrule
\end{tabularx}
\end{table*}

\subsection{Transparency Dashboard}
\begin{table*}[h!]
\centering
\renewcommand{\arraystretch}{1.25}
\footnotesize
\begin{tabularx}{\textwidth}{@{}p{4.2cm}Xp{3.6cm}@{}}
\toprule
\textbf{Prompt (intent + verbatim)} & \textbf{Reply (gpt-5-chat-latest, best-of-5, truncated)} & \textbf{What it demonstrates} \\
\midrule
\textbf{Procedural}\newline why does the dashboard say 'approximately' next to the numbers? & It says "approximately" because the dashboard adds a small amount of random noise to each number before publishing. That's called *Laplace noise*. \ldots & \textit{Engages with the user's situation; no refusal-bot tone.} \\
\textbf{Practical worry}\newline i'm a contractor, not a regular employee. can i still file? & Yes, you can. Section 2(f) of the POSH Act covers not just employees but also contractors, interns, visitors, and clients. You have the same right to file a complaint under Section 9 as anyone on the company's rolls. & \textit{Procedural answer grounded in the right section / rule.} \\
\textbf{Boundary}\newline my colleague did something to me — was that legally sexual harassment? & I'm really sorry that something happened that's worrying you. I can't decide whether what occurred legally counts as sexual harassment — that's for the Internal Complaints Committee (ICC) to determine. \ldots & \textit{Role boundary holds; reply redirects without lecturing.} \\
\textbf{Decision-help}\newline will my manager retaliate against me if i file? what protections exist? & The law says retaliation is not allowed. Under Section 19(i) of the POSH Act, sexual harassment is treated as misconduct under service rules, and any retaliation—like bad appraisals, transfers, or threats for filing—is itself a violation. \ldots & \textit{Procedural answer grounded in the right section / rule.} \\
\textbf{Edge case}\newline what actually happens after i file? how long does it take? & Once you file under Section 9, the Internal Complaints Committee (ICC) first checks that the complaint is within its scope and time window. Then: 1. **Acknowledgment:** You'll get a written acknowledgment. 2. \ldots & \textit{Procedural answer grounded in the right section / rule.} \\
\bottomrule
\end{tabularx}
\end{table*}
